\chapter{Sécurité}\label{chap:securite}

\section{Ce qui est repris, ce qui est adapté}
Les mécanismes de l'atelier~2 sont conservés ; leur emplacement change, parce que les servlets ont disparu.
Le tableau~\ref{tab:menaces} relie chaque menace à sa parade dans la version JSF.

\begin{table}[htbp]
\centering
\caption{Menaces et parades}\label{tab:menaces}
\small
\begin{tabularx}{\linewidth}{@{}p{3.2cm}X@{}}
\toprule
\textbf{Menace} & \textbf{Parade}\\
\midrule
Injection SQL & Requêtes JPQL paramétrées, aucun assemblage de texte (inchangé).\\
XSS & Facelets échappe par défaut toute expression \code{\#\{...\}} ; aucune sortie en \code{escape="false"} ; CSP \code{script-src 'self'}.\\
CSRF & Jeton \code{\_csrf} par session, inséré par \code{t:csrf}, vérifié en temps constant par \code{SecurityFilter} sur tout POST. À cela s'ajoute le \code{ViewState} de JSF.\\
Fixation de session & \code{changeSessionId()} et renouvellement du jeton à la connexion (\code{UserSession.ouvrir}).\\
Force brute & \code{LoginGuard} : 5 essais par 15 minutes, par couple IP / adresse.\\
Redirection ouverte & \code{Chemins.cheminSur} n'accepte que les chemins internes (refus de \code{//}, d'URL absolue, d'antislash, de saut de ligne).\\
Contrôle d'accès & \code{AuthFilter} par motif d'adresse : \code{/admin/*} exige le rôle \code{ADMIN}, \code{/panier.xhtml} le rôle \code{CLIENT} (réponse 403 sinon).\\
Mots de passe & PBKDF2 avec sel aléatoire ; jamais renvoyés à la page.\\
Accès aux composants & Contrainte de sécurité sur \code{/resources/components/*}.\\
\bottomrule
\end{tabularx}
\end{table}

\section{Filtres}
Deux filtres s'appliquent, dans cet ordre déclaré dans \code{web.xml}.

\textbf{\code{SecurityFilter}} (toutes les adresses) : fixe l'encodage UTF-8, ajoute les en-têtes de sécurité
(\code{Content-Security-Policy}, \code{X-Content-Type-Options}, \code{X-Frame-Options}, \code{Referrer-Policy}) et \code{Cache-Control: no-store}, génère le jeton CSRF à la
première visite et refuse tout \code{POST} dont le jeton est absent ou différent.

\textbf{\code{AuthFilter}} (\code{/panier.xhtml}, \code{/admin/*}) : sans utilisateur en session, redirige vers
\code{connexion.xhtml?requis=1} et mémorise l'adresse demandée. Après connexion, l'utilisateur retourne où il voulait aller ; cette adresse
est revérifiée par \code{Chemins.cheminSur} avant la redirection.

\section{Politique de sécurité du contenu}
\begin{lstlisting}[language=Java,caption={Politique appliquée à toutes les réponses},label={lst:csp}]
default-src 'self'; img-src 'self' data: https:;
style-src 'self' 'unsafe-inline'; script-src 'self';
frame-ancestors 'none'; form-action 'self'; base-uri 'self'
\end{lstlisting}

\code{script-src 'self'} est le c\oe{}ur : seuls les scripts servis par l'application s'exécutent. \code{frame-ancestors 'none'}
interdit l'inclusion de l'application dans un cadre d'un autre site (détournement de clic) ; \code{form-action 'self'} empêche un
formulaire injecté d'envoyer ses données ailleurs. \code{style-src} garde \code{'unsafe-inline'} parce que JSF et les
composants produisent des attributs \code{style} ; c'est un compromis assumé, le risque d'un style injecté étant bien
moindre que celui d'un script.

\section{Contrôle d'accès}
Un utilisateur qui n'a pas le rôle requis (par exemple un client devant une page d'administration) reçoit une réponse 403 affichée dans la page d'erreur de
l'application, pas la page technique de Tomcat.

\fig{16-acces-refuse}{Accès refusé : un client demande une page d'administration}{fig:403}
